O objetivo deste trabalho foi desenvolver e avaliar uma ferramenta computacional para o dimensionamento e otimização de fundações superficiais do tipo sapata isolada, integrando critérios estruturais, geotécnicos e geométricos em um único ambiente de projeto. Para tal, foi proposta uma formulação de otimização mono-objetivo baseada na minimização do volume de concreto, empregando restrições normativas e construtivas, resolvida por meio do método \textit{Efficient Global Optimization} (EGO).

Com base nos resultados obtidos, destacam-se as seguintes conclusões:
\begin{enumerate}
    \item A formulação do problema de dimensionamento como um problema de otimização com restrições mostrou-se adequada para representar, de forma concomitante, os  critérios normativos aplicados às sapatas isoladas, incluindo verificações de tensões no solo, punção, compatibilidade geométrica entre o pilar e a sapata, além da sobreposição entre elementos de fundação.
    \item O EGO mostrou-se eficiente para problemas com dimensionalidade crescente, como no dimensionamento de múltiplas sapatas, onde cada elemento contribui com três variáveis ($h_x$, $h_y$, $h_z$).
    \item A abordagem baseada em modelo substituto garante uma busca global eficaz, como evidenciado pela superioridade frente a amostragem aleatória (Monte Carlo) em termos de qualidade da solução final, validando sua capacidade de evitar convergência prematura para ótimos locais.
    \item Os exemplos analisados demonstraram que a metodologia proposta é capaz de convergir para soluções geometricamente proporcionais e tecnicamente viáveis, atendendo simultaneamente às restrições estruturais, geotécnicas e geométricas, além de apresentar uma média de redução de 40,79\%, até 85,87\%, do consumo de concreto quando comparada a abordagem MMC que representa a técnica de escritório baseada em tentativas sucessivas.
    \item A verificação explícita de sobreposição entre sapatas mostrou-se um diferencial relevante do modelo, permitindo identificar interferências geométricas em planta e incorporá-las diretamente ao processo de otimização por meio das funções de restrição normalizadas.
    \item Nos casos em que os elementos de fundação apresentam elevado grau de proximidade, os resultados indicaram que a estratégia adotada, baseada apenas na variação das dimensões das sapatas, pode conduzir a soluções que não convergem para a solução ótima do problema. Isso evidencia que a metodologia, em sua forma atual, é mais adequada a arranjos de fundações sem sobreposição severa. Dessa forma, a sobreposição poderia ser resolvida com a implementação de uma segunda etapa de otimização, na qual a associação das sapatas deve ser verificada.
    \item A implementação da metodologia em uma plataforma computacional com interface \textit{web} demonstrou a viabilidade prática da abordagem, permitindo a aplicação direta do modelo em problemas reais de projeto, com entrada de dados via planilha e exportação automatizada dos resultados de dimensionamento e verificação.
    \item Em casos com mais de um elemento de fundação, a otimização global do problema não exige que todas as fundações apresentem restrições ativas. Sob a perspectiva de projeto individual, é desejável que cada elemento opere próximo a pelo menos um estado limite, garantindo o aproveitamento eficiente de material e racionalidade construtiva.
   
\end{enumerate}

De modo geral, a metodologia proposta mostrou-se uma ferramenta eficiente e promissora para o projeto racional e automatizado de sapatas isoladas, contribuindo para a padronização do processo de dimensionamento, a redução da subjetividade do projetista e o uso mais eficiente de materiais. Como trabalhos futuros, recomenda-se a incorporação de abordagens probabilísticas e de análise de confiabilidade, a consideração de sapatas associadas em situações de elevada proximidade entre pilares e a ampliação do modelo para incluir efeitos de recalque, interação solo-estrutura e detalhamento estrutural, de modo a expandir o campo de aplicação da ferramenta desenvolvida.
    







% The aim of this paper was to evaluate the reliability index of reinforced concrete beams affected by carbonation-induced corrosion. Reinforced concrete sections with a compressive strength of 25 MPa, subjected to a class C4 environment, were analyzed. The principal conclusions are as follows:
% \begin{enumerate}
%     \item In this environmental condition, the depassivation of the reinforcement due to carbonation occurred after a period of 39 years, with a standard deviation of 18 years. This suggests that, adopting a 95\% confidence interval, the carbonation corrosion process can potentially initiate after only 3.7 years.
%     \item The longitudinal reinforcement rate influenced the period in which the target reliability index of 3.8 proposed by the JCCS [36] and the Model Code 2010 [32] for the Ultimate Limit State was reached. Adopting a concrete cover of 2.5 cm, beams with a longitudinal reinforcement ratio of 0.2\% reached this target reliability index after 60 years. However, in beams with a longitudinal reinforcement ratio of 0.9\%, this target reliability index was reached after a period of only 30 years.
%     \item A 12\% variation in the target reliability index at the Ultimate Limit State after 50 years was observed when the concrete cover was varied between 2 and 3 cm. A concrete cover of only 2 cm can result in an average reliability index of 3.28, or 2.77 if a 95\% confidence level is adopted. On the other hand, a concrete cover of 3 cm yields an average reliability index of 3.66, or up to 3.25 if a 95\% confidence level is assumed. This demonstrates the importance of increasing the concrete cover to achieve higher reliability indices. A concrete cover value of 3 cm would appear to be sufficient for structures subjected to environmental exposure class C4 to ensure that the target reliability index of 3.8 is met for the majority of structures.
%     \item Considering 3 cm of concrete cover and age over 75 years, the average reliability index value was below the minimum limit of 3.1 recommended in the design standards for the Ultimate Limit State. This indicates that an increase in the longitudinal reinforcement cover is necessary to ensure that the minimum value of 3.1 is met after 50 years.
%     \item \textcolor{red}{Based on the significant correlation depicted in Figure 5.3, a predictive model was developed to assess the structural reliability of concrete beams subjected to corrosive conditions. The resulting expression, presented as Equation (5.1), constitutes a reliable tool for evaluating the safety of concrete elements in corrosion-prone environments and is suitable for rapid condition assessments during field inspections.}
%     \item \textcolor{red}{The technique proposed in this paper has proven to be an effective method for assessing the reliability index in the Ultimate Limit State, making it a valuable tool for evaluating the long-term maintenance needs of structures. It enables the estimation of the initiation time of corrosion due to carbonation and provides the corresponding reliability index, helping in decision-making regarding when interventions are necessary to restore the target safety levels specified in design guidelines. The results indicate a reduction in the reliability index over time, especially within the first 50 years, highlighting the importance of timely maintenance to maintain structural safety. Moreover, considering the growing impact of extreme climatic conditions, there is an increasing need to reassess the partial coefficients used in design. These coefficients may need adjustment to better reflect the accelerated degradation caused by changing environmental factors, ensuring that structural safety is accurately accounted for in the face of evolving climate conditions.}
%     \item \textcolor{red}{The proposed model presents some simplifications that should be acknowledged. In particular, the analysis does not consider the reduction of steel strength caused by corrosion, nor the potential loss of concrete cross-section due to spalling as corrosion progresses. These assumptions may lead to conservative or optimistic outcomes depending on the specific structural and environmental conditions.}
%     \item \textcolor{red}{In addition, the results are valid for exposed concrete structures (without painting or protective coating), under the same environmental exposure conditions considered in the analysis. Therefore, the findings should not be generalized to different exposure scenarios without further adjustments.}    
    
% \end{enumerate}

% \textcolor{red}{The proposed methodology offers valuable insights for concrete structures and supports decision-making processes for the maintenance and life-cycle management of civil infrastructure. Furthermore, it is instrumental in prioritizing structures for detailed investigation and intervention, particularly within large asset inventories, thereby facilitating strategic resource allocation for optimized maintenance planning.} However, it should be noted that the results presented in this paper were obtained based on the humidity and temperature conditions presented in the methodology. They may differ under other climatic conditions, and the proposed predictive equation should be adjusted accordingly.


